% !TEX root = ../main.tex
\section{Conclusion}
\label{sec:conclusion}

I have argued for two connected theses.

First, the assertion operator of a language model, $\Out_M$, is hyperintensional
and non-normal. It fails K, M, D, 4, and 5, and critically RE, as a constitutive
matter of an architecture that computes over formulations rather than
propositions. That its outputs are not \emph{knowledge} follows from factivity
alone and needs none of this apparatus. What the apparatus establishes is the
doxastic claim: since belief in the received formal tradition is a normal
operator, and $\Out_M$ is not even classical, these systems are not doxastic
systems in that sense. Their assertion behaviour admits no Kripke semantics for
any accessibility relation, and no bare neighborhood semantics either, since
every neighborhood frame validates RE; what it requires is a hyperintensionalised
neighborhood model with neighborhoods over formulations, or an impossible-worlds
semantics. The system-level conclusion rests on a premise about internal states,
and I have argued for the modest version of it: not that no prompt-invariant
truth-correlated state exists---the evidence suggests one does---but that no such
state governs what the system asserts. The dissociation between what these
systems represent and what they say is not an embarrassment for the account. It is
the account's central evidence.

The resulting characterisation is an extended, monadic descendant of Stich's
sub-doxastic category. Extended twice over: because these systems fail the
inferential-isolation mark in an unexpected direction, exhibiting promiscuity
without discipline rather than encapsulation; and because there is no
personal-level believer for the states to be sub-personal relative to. Both are
limits of the analogy I have stated rather than elided, and what survives them is
a monadic core---informational, behaviour-guiding, closure-free---that names a
kind of system cognitive science never needed a word for, because in creatures
such states always come embedded in a believer.

Second, once grounding is prised apart from truth and indexed to the context of
use, hallucination is properly defined as ungrounded assertion, with no
truth-value condition. Faithful transmission of a bad source is error. Assertion
grounded in evidence that is stale or of the wrong jurisdiction is neither error
nor fabrication but a third thing with its own remedy. And a lucky truth
generated without evidential connection is a hallucination despite being
true---not a Gettier case, since no justification is present to be undermined, but
a failure of safety, an assertion that would have been produced whatever the fact
had been. The Gettier structure appears one step downstream, in the user who
takes fluency for justification and thereby acquires a justified true belief that
is not knowledge. Truth-matching benchmarks misclassify both off-diagonal cells,
and Proposition~\ref{prop:independence} shows the misclassification is
unbounded: accuracy constrains the grounding rate not at all. The correct metric
is the rate at which assertion outruns applicable evidence, measured by
attribution audit---and measurable only for systems built so that the evidence is
identified at the point of assertion.

The two theses are one diagnosis viewed from two sides. Hallucination in the
defined sense is the behavioural expression of hyperintensionality: where
assertion probability is fixed by surface statistics, assertion floats free of
evidence, and sometimes the float lands on the truth. Understanding these systems
epistemically is therefore not a matter of deciding how much they know. It is a
matter of theorising, with the right formal tools, a genuinely new kind of
system---informational without being doxastic, assessable without being
appraisable, and reliable, when it is reliable, in a way that has to be
engineered and measured rather than assumed. Free-standing sub-doxastic engines,
whose epistemology begins where the epistemology of belief leaves off.
